\documentclass[11pt]{article}
\usepackage[margin=1.15in]{geometry}
\usepackage{amsmath,amssymb}
\usepackage{graphicx}
\usepackage{booktabs}
\usepackage[hidelinks]{hyperref}
\usepackage{microtype}
\graphicspath{{../figures/}}

\newcommand{\fig}[3]{\begin{figure}[ht]\centering
\includegraphics[width=#2\textwidth]{#1}\caption{#3}\end{figure}}

\title{How the Primes Work, in Five Pictures}
\author{VFD Programme\thanks{Companion to \emph{The Prime Field Through One
Instrument} (the verification ledger) in this repository:
\texttt{github.com/vfd-org/closure-positivity-lab}. That paper verifies;
this one explains. Every figure here is computed, not drawn --- the
generating script is \texttt{figures/make\_pictures.py}, and two of the
pictures are interactive at \texttt{explorables/}.}}
\date{June 2026}

\begin{document}
\maketitle

\begin{abstract}\noindent
The ledger paper in this repository reports that ten prime-distribution
phenomena pass numerical gates; it does not explain \emph{why}. This
companion supplies the why, at the level where the why is actually known:
in five pictures. (1)~Patterns of primes live or die cell-by-cell on each
prime's \emph{wheel}, and the surviving fractions multiply into every
density law Hardy and Littlewood wrote down. (2)~That one multiplication
predicts twin counts to $0.003\%$, the doubled abundance of gap-6 pairs,
and Goldbach counts per~$N$. (3)~Read through the explicit formula, the
primes are a sum of waves whose spectrum peaks exactly at the zeros of the
zeta function --- we plot this, and the peaks land on independently
computed zeros. (4)~A one-dimensional golden-ratio quasicrystal shows the
same signature --- sharp spectrum, no period --- which is precisely Dyson's
reason for suggesting that the Riemann Hypothesis is secretly a statement
about quasicrystals; this is where the question becomes geometric.
(5)~Everything open in the subject is one wall: the correlations of those
zeros. All mathematics here is classical and attributed; the pictures and
their code are ours. Nothing below proves, or claims to prove, any open
conjecture.
\end{abstract}

\section{The wheels: where the densities come from}

Look at the whole number line through a wheel with three cells --- the
residues $0,1,2$ modulo $3$. Every integer lands on one cell; primes
(other than $3$ itself) only ever land on cells $1$ and $2$, because cell
$0$ means ``divisible by $3$.''

Now place a \emph{pattern} on the wheel instead of a single number: the
twin pattern $(n,\,n{+}2)$. If $n$ is on cell $0$, the pattern dies ($n$
is divisible by 3). If $n$ is on cell $1$, then $n+2$ is on cell $0$:
dead again. Only $n$ on cell $2$ survives. \textbf{One cell of three.}

Compare with what independence would give: two unrelated numbers each
avoid cell $0$ with probability $2/3$, so independent survival would be
$(2/3)^2 = 4/9$. The twin pattern's actual survival is $1/3 = 3/9$. The
wheel taxes twins by a factor
\[
\frac{1/3}{(2/3)^2} \;=\; \frac{3}{4}
\;=\;\frac{p(p-2)}{(p-1)^2}\Big|_{p=3},
\]
and the same computation on every prime wheel $p>2$ gives the factor
$p(p-2)/(p-1)^2$: the pattern occupies two cells, so it dies on two of
$p$, surviving $p-2$, against an independence baseline of $((p-1)/p)^2$.
Multiply over all odd wheels:
\[
\prod_{p>2}\frac{p(p-2)}{(p-1)^2}
\;=\;\prod_{p>2}\Bigl(1-\frac{1}{(p-1)^2}\Bigr)
\;=\;C_2\;=\;0.6601618\ldots
\]
This is the twin-prime constant of Hardy and Littlewood \cite{HL1923},
and now you know what it \emph{is}: the total tax the wheels charge the
twin pattern, relative to independence. The predicted count below $x$ is
$2C_2\int_2^x dt/\log^2 t$, and the sieve agrees to $0.003\%$ at
$x=5\times10^7$ (ledger row~1).

\fig{fig_wheels.png}{0.82}{The mod-3 and mod-5 wheels. Green cells:
starting positions where the pattern avoids every forbidden cell. Twins
survive 1 cell of 3 (top left); the gap-6 pattern survives 2 of 3 (top
right) --- that single picture is the whole reason gap-6 pairs are twice
as common as twins. Bottom right: $(n,n{+}2,n{+}4)$ kills every cell mod
3, so only one such triple exists in all of mathematics --- $(3,5,7)$,
where the dead cell is the prime 3 itself.}

Two more readings of the same picture. \textbf{Why gap 6 doubles:} the
pattern $(n,\,n{+}6)$ doesn't move on the wheel of 3 (adding 6 is adding
0), so it dies only on cell $0$: two cells of three survive ---
\emph{twice} the twin survival, and no other wheel distinguishes the two
patterns. Hence gap-6 pairs are exactly twice as abundant; the sieve says
$1.9950$. \textbf{Why $(n,n{+}2,n{+}4)$ is impossible:} on the wheel of
3 the three offsets cover all three cells, so every starting position
dies. The only escape is the pattern that \emph{includes} 3 itself:
$(3,5,7)$, once, never again. The ledger checks the analogous exact fact
on the wheel of 5: the only twin pair $(p,p{+}2)$ with $p\equiv3\pmod 5$
is $(3,5)$, because cell 3 forces $5\mid p+2$.

In the programme's language each wheel is a \emph{closure condition at the
place $p$}: a finite cyclic screen that a pattern either clears or does
not. The density of any pattern is its product of survival ratios over
all screens --- an \emph{all-finite-places product}, the same shape as the
total-positivity law studied elsewhere in this repository. That parallel
is structural, not decorative: in both cases the decidable content lives
place-by-place, and the open content lives in how infinitely many places
conspire.

\section{One multiplication, every pattern}

The wheel computation generalises mechanically: for a pattern of $k$
offsets, count surviving cells $w(p)$ on each wheel, divide by the
independence baseline, multiply. The result (the \emph{singular series})
predicts, with \emph{no further input}:

\begin{itemize}
\item the relative abundance of every even gap (figure below: 105 gaps,
  predictions as dots, sieve counts as bars --- worst disagreement
  $0.25\%$);
\item Goldbach counts per individual $N$ --- the wheels of the primes
  \emph{dividing} $N$ relax their conditions, so $N$ divisible by 3 has
  about twice the representations of its neighbours (verified: 200
  consecutive even $N$ near $10^6$, mean deviation $0.65\%$);
\item the twin count itself, to $0.003\%$.
\end{itemize}

\fig{fig_gap_bullseyes.png}{0.95}{All even gaps to 210: measured abundance
(bars) against the wheel product (dots). The spikes are gaps divisible by
3, 15, 105\ldots\ --- more surviving cells, more pairs. Interactive
version: \texttt{explorables/gap-slider.html}.}

What the wheels do \emph{not} give is infinitude: they say twin pairs
keep their density \emph{if anything like statistical regularity holds},
but cannot rule out the counts simply stopping. That question changes
character entirely, and needs the second half of the story.

\section{The duality: primes are waves, zeros are the spectrum}

Riemann's explicit formula is an exact identity, and it has a physical
reading. Place every prime $p$ on a logarithmic line and let it broadcast
a wave of frequency $\log p$ (with weight $\log p/\sqrt p$, and its powers
chiming in quietly). Sum the waves and take the spectrum --- the same
operation a crystallographer performs when she shines X-rays through a
crystal and reads the diffraction pattern.

\fig{fig_prime_diffraction.png}{0.95}{The spectrum of the prime waves
(orange), computed from the primes below $10^6$ with a smooth taper, and
\emph{nothing else}; dashed lines are the first zeta zeros, computed
independently in PARI/GP from the zeta function. The peaks land on the
zeros: 13 of the first 15 align within $0.15$ (the misses are close
zero-pairs merged at this resolution).}

The picture says what the explicit formula says: \textbf{the zeros of
$\zeta$ are the spectrum of the primes}. The identity runs both ways ---
zeros reconstruct prime counts, primes reconstruct zeros (the companion
paper's prime-wave reconstruction of the critical line is this same
duality run in reverse for our icosian $L$-function). Every question the
wheels cannot answer is a question about this spectrum: the \emph{error}
in the prime count is a sum over zeros; the Chebyshev bias is an
interference effect (the prime-\emph{square} wave loads one residue
class); twin-prime infinitude is governed by correlations \emph{between}
zeros. The wheels are the local physics; the spectrum is the global
physics; the explicit formula is the exchange rate.

\section{The crystal question: where it becomes geometry}

A crystal produces sharp diffraction peaks because it is periodic. In
1982 Shechtman saw sharp peaks with a symmetry no period allows; the
resolution is the \emph{quasicrystal}, and its one-dimensional archetype
is the Fibonacci chain --- two tile lengths in ratio $\varphi$ (the golden
ratio), arranged by the substitution $A\to AB$, $B\to A$, never
repeating.

\fig{fig_fibonacci.png}{0.95}{Computed diffraction of 3{,}000 points.
Top: the Fibonacci $\varphi$-chain --- razor-sharp peaks, no period.
Bottom: the same number of points placed at random --- featureless. Sharp
spectrum without periodicity is possible, and golden-ratio geometry is
its canonical home.}

Now hold figure~3 next to figure~4's top panel. The primes (suitably
weighted, on the log line) also show a sharp spectrum without any period
--- \emph{provided the peaks all sit where we plotted them}, namely on
one vertical line in the complex plane. That proviso is the Riemann
Hypothesis. This observation is due to Dyson \cite{Dyson}, who proposed
classifying one-dimensional quasicrystals as a route to RH: if the primes'
spectrum is genuinely point-like, the primes are (in this precise sense)
a quasicrystal, and quasicrystals might be classifiable.

We state plainly what this is and is not. It is the one place where the
deepest open question about primes becomes a question of geometry --- the
geometry of aperiodic order, whose golden-ratio examples (Fibonacci
chains, Penrose tilings, the icosians of this repository) are the native
objects of our wider programme. It is \emph{not} a result: nobody has the
classification, we do not claim it, and the honest status of every
approach we have tested is recorded in the companion papers --- including
the negative results.

\section{The wall, and the map}

Assemble the picture. Everything decidable about prime patterns we have
met lives \emph{locally}: wheels, survival counts, one multiplication ---
verified across ten phenomena in the ledger at the $0.003\%$--$0.25\%$
level, with three theorem-grade anchor rows calibrating the instrument.
Everything open lives \emph{globally}, and it is one object, not many:
the fine structure of the zero spectrum. Twin primes, Polignac, Goldbach,
Cram\'er's gap bound, the unconditional Chebyshev density --- each is a
different question about the same correlations (ledger row~10 measures
them directly: spacing statistics of $1{,}983$ computed zeros match the
random-matrix law within the documented finite-height corrections, with
level repulsion suppressed $70\times$ relative to Poisson).

\begin{center}\small
\begin{tabular}{@{}lll@{}}
\toprule
 & the local layer (wheels) & the global layer (spectrum) \\
\midrule
what it fixes & whether a pattern occurs; its density & infinitude;
error terms; extremes \\
status & decidable, computed, verified & open in every conjectural case \\
the ``how'' & counting cells on finite screens & unknown; best picture:
diffraction \\
where geometry enters & each place is a finite cyclic screen &
Dyson's quasicrystal question \\
\bottomrule
\end{tabular}
\end{center}

That is the full picture of the primes as honestly as we can paint it:
a local mechanism that is completely understood and astonishingly
accurate, a global mechanism that is one open question wearing many
costumes, an exact dictionary between them, and a geometric reading of
the open half that remains a question. The ledger paper carries the
gates, artifacts and falsifiers behind every number quoted here; the
explorables let you drive the two halves yourself.

\begin{thebibliography}{9}\small
\bibitem{HL1923} G.~H. Hardy, J.~E. Littlewood, \emph{Some problems of
`Partitio Numerorum' III: on the expression of a number as a sum of
primes}, Acta Math.\ 44 (1923) 1--70.
\bibitem{Dyson} F.~Dyson, \emph{Birds and frogs}, Notices Amer.\ Math.\
Soc.\ 56 (2009) 212--223 (the quasicrystal proposal, p.~220).
\bibitem{BaakeGrimm} M.~Baake, U.~Grimm, \emph{Aperiodic Order, Vol.~1},
Cambridge Univ.\ Press (2013).
\bibitem{Montgomery} H.~Montgomery, \emph{The pair correlation of zeros
of the zeta function}, Proc.\ Sympos.\ Pure Math.\ 24 (1973) 181--193.
\bibitem{Odlyzko} A.~Odlyzko, \emph{On the distribution of spacings
between zeros of the zeta function}, Math.\ Comp.\ 48 (1987) 273--308.
\bibitem{RubinsteinSarnak} M.~Rubinstein, P.~Sarnak, \emph{Chebyshev's
bias}, Experiment.\ Math.\ 3 (1994) 173--197.
\bibitem{Maynard} J.~Maynard, \emph{Small gaps between primes}, Ann.\ of
Math.\ 181 (2015) 383--413.
\bibitem{Ledger} VFD Programme, \emph{The Prime Field Through One
Instrument: a ten-row phenomena ledger}, this repository (2026).
\bibitem{Diffraction} VFD Programme, \emph{Closure diffraction and the
icosian $L$-function}, this repository (2026).
\end{thebibliography}

\end{document}
